% 保持共同预测不变，仅比较行动与后果；节点少，按简单框图用TikZ。
\begin{tikzpicture}[figure picture,foundation picture,
 trust node/.style={draw=FigureStroke,line width=.7pt,rounded corners=1.5pt,align=center,minimum height=13mm,inner sep=2mm},
 trust consequence/.style={trust node,minimum height=24mm,font=\FigureNoteFont\fontsize{7.5}{9.5}\selectfont}]
 \path[use as bounding box] (0,0) rectangle (163.2,65);
 \node[trust node,fill=FigureBlueFill,minimum width=23mm] (input) at (12,32) {正常邮件};
 \node[trust node,fill=MutedPurpleFill,minimum width=28mm] (model) at (52,32) {相同预测\\$p=f(\bm x)$};
 \node[trust node,fill=MutedGreenFill,minimum width=30mm] (archive) at (98,49) {隔离\\保留原件};
 \node[trust node,fill=FigureRedFill,minimum width=30mm] (delete) at (98,15) {永久删除\\删除原件};
 \node[trust consequence,fill=MutedGreenFill,minimum width=39mm] (recover) at (143,49) {可能找回原件\\承担找回成本\\未必及时找回};
 \node[trust consequence,fill=FigureRedFill,minimum width=39mm] (lost) at (143,15) {原件无法恢复\\可能错过期限};
 \draw[foundation flow] (input.east)--(model.west);
 \draw[foundation flow] (model.east)--(75,32)|-(archive.west);
 \draw[foundation flow] (75,32)|-(delete.west);
 \node[foundation edge label] at (70.5,37) {误拦};
 \draw[foundation flow] (archive.east)--(recover.west);
 \draw[foundation flow] (delete.east)--(lost.west);
\end{tikzpicture}
